% !TEX root = ../main.tex
%% Corrigés — Chapitre 5 — Logarithme népérien
%% Tous les résultats ont été vérifiés par calcul formel (sympy).

\corrigesChapitre{5}{Logarithme népérien}

\begin{deuxcolonnes}

\begin{corrige}{1}{Propriétés algébriques}
\begin{questions}
  \item $\ln(8\times2)=\ln16=\res{4\ln2}$
  \item $\ln\left(\frac{27}{3}\right)=\ln9=\res{2\ln3}$
  \item $\ln25+\ln8=\res{\ln200}$
  \item $\ln\left(\frac{16}{9}\times\frac34\right)=\res{\ln\frac43}$
  \item $3-2-1=\res{0}$
\end{questions}
\end{corrige}

\begin{corrige}{2}{Puissances et racines}
\begin{questions}
  \item $\ln\left(e^{1/2}\right)=\res{\frac12}$
  \item $\res{5}$
  \item $\ln\left(e^{-2}\right)=\res{-2}$
  \item $\ln\left(e^{1/3}\right)=\res{\frac13}$
  \item $2\times3-1=\res{5}$
\end{questions}
\end{corrige}

\begin{corrige}{3}{Dérivées}
\methode{$\big(\ln u\big)'=\frac{u'}{u}$}
\begin{questions}
  \item $\res{f'(x)=\dfrac{3}{3x+1}}$ sur $\intOO{-\frac13}{+\infty}$
  \item $\res{g'(x)=\dfrac{2x}{x^{2}-4}}$
  \item $h'(x)=2x\ln x+x^{2}\times\frac1x$\\ $\res{h'(x)=x(2\ln x+1)}$ sur $\intOO{0}{+\infty}$
  \item $\res{k'(x)=\dfrac{1-\ln x}{x^{2}}}$
  \item $m(x)=\frac12\ln\left(x^{2}+1\right)$ : $\res{m'(x)=\dfrac{x}{x^{2}+1}}$ sur $\R$
  \item $\res{p'(x)=\dfrac{2x}{x^{2}+4}}$ sur $\R$
  \item $3x^{2}+2x-1=(3x-1)(x+1)>0$ sur $\intOO{-\infty}{-1}\cup\intOO{\frac13}{+\infty}$ : $\res{q'(x)=\dfrac{6x+2}{3x^{2}+2x-1}}$
\end{questions}
$p'(x)$ est du signe de $x$ : $p$ décroît sur $\intOF{-\infty}{0}$, croît sur $\intFO{0}{+\infty}$ ; minimum $p(0)=\ln4$.
\end{corrige}

\begin{corrige}{4}{Équations}
\begin{questions}
  \item Il faut $x>\frac13$ ; $x+2=3x-1$ donne $\res{S=\ens{\frac32}}$.
  \item $\res{S=\ens{e^{2}}}$
  \item Il faut $x>2$. Pour $x>2$, $\frac{x^{2}-4}{x-2}=x+2$, et $\ln(x+2)=1$ donne $x=e-2\approx0{,}72$, qui ne convient pas : $\res{S=\varnothing}$.
  \item Il faut $x>1$ ; $x^{2}-1=8$ donne $x=3$ : $\res{S=\ens{3}}$.
  \item Il faut $x>0$ ; $x^{2}=x^{2}+3$ est impossible : $\res{S=\varnothing}$.
\end{questions}
\end{corrige}

\begin{corrige}{5}{Inéquations}
\begin{questions}
  \item Il faut $x>0$ ; $x+1<2x\iff x>1$ : $\res{\intOO{1}{+\infty}}$
  \item $x\geq e^{-1}$ : $\res{\intFO{\frac1e}{+\infty}}$
  \item $0<2x+3\leq1$ : $\res{\intOF{-\frac32}{-1}}$
  \item $x^{2}-1>3\iff x^{2}>4$ : $\res{\intOO{-\infty}{-2}\cup\intOO{2}{+\infty}}$
  \item Il faut $x>2$ ; $x(x-2)\leq8\iff(x-4)(x+2)\leq0$ : $\res{\intOF{2}{4}}$
\end{questions}
\end{corrige}

\begin{corrige}{6}{Limites et croissances comparées}
\begin{questions}
  \item Croissances comparées : $\res{0}$
  \item Croissances comparées : $\res{0}$
  \item Nombre dérivé de $x\mapsto\ln(1+x)$ en $0$ : $\res{1}$
  \item Inverse de $\frac{\ln x}{x}\to0^{+}$ : $\res{+\infty}$
  \item Nombre dérivé de $\ln$ en $1$ : $\frac11=\res{1}$
  \item $\frac{2\ln x}{x}\to\res{0}$
\end{questions}
\end{corrige}

\begin{corrige}{7}{Études de fonctions}
\begin{questions}
  \item Sur $\intOO{\frac12}{+\infty}$, $f'(x)=\frac{2}{2x-1}>0$ : $f$ croît de $-\infty$ à $+\infty$ ; asymptote $x=\frac12$.
  \item Sur $\intOO{0}{+\infty}$, $g'(x)=\ln x+1$ : $g$ décroît sur $\intOF{0}{\frac1e}$, croît ensuite ; $\res{\min g=-\frac1e}$ ;
    $\lim_{0^{+}}g=0$, $\lim_{+\infty}g=+\infty$.
  \item Sur $\intOO{-\infty}{1}\cup\intOO{2}{+\infty}$, $h'(x)=\frac{2x-3}{x^{2}-3x+2}$ : $h$ décroît sur $\intOO{-\infty}{1}$, croît sur
    $\intOO{2}{+\infty}$ ; limites $+\infty$ en $\pm\infty$ et $-\infty$ en $1^{-}$ et $2^{+}$.
  \item Sur $\intOO{0}{+\infty}$, $k'(x)=\frac{1-2\ln x}{x^{3}}$ : $k$ croît sur $\intOF{0}{\sqrt e}$, décroît ensuite ;
    $\res{\max k=\frac{1}{2e}}$ ; $\lim_{0^{+}}k=-\infty$, $\lim_{+\infty}k=0$.
  \item Sur $\intOO{-\infty}{-1}\cup\intOO{1}{+\infty}$, $m'(x)=-\frac{2}{x^{2}-1}<0$ : $m$ décroît sur chaque intervalle ;
    limite $0$ en $\pm\infty$, $-\infty$ en $-1^{-}$, $+\infty$ en $1^{+}$.
\end{questions}
\end{corrige}

\end{deuxcolonnes}

%% ---------------------------------------------------------------------
\begin{corrige}{8}{Modélisation}
\begin{questions}
  \item $e^{3k}=2$ donne $\res{k=\frac{\ln2}{3}}$. Puis $e^{kt}=100$ donne $t=\dfrac{\ln100}{k}=\dfrac{3\ln100}{\ln2}\approx19{,}9$ : environ $\res{20\text{ heures}}$.
  \item $1{,}05^{n}\geq2\iff n\geq\dfrac{\ln2}{\ln1{,}05}\approx14{,}2$. Les intérêts étant versés chaque année, le capital double
    au bout de $\res{15\text{ ans}}$ ($1{,}05^{14}\approx1{,}980$ et $1{,}05^{15}\approx2{,}079$).
  \item $80e^{-10k}=40\iff e^{-10k}=\frac12\iff\res{k=\frac{\ln2}{10}\approx0{,}069}$.
  \item $e^{-30\lambda}=\frac12$ donne $\res{\lambda=\frac{\ln2}{30}}$. Proportion restante : $e^{-100\lambda}=2^{-10/3}\approx\res{9{,}9\,\%}$.
  \item $\ln\left(\frac{I}{I_0}\right)=8$, donc $\res{I=e^{8}I_0\approx2\,981\,I_0}$.
\end{questions}
\end{corrige}

%% ---------------------------------------------------------------------
\begin{corrige}{9}{Synthèse 1}
\begin{questions}
  \item $x^{2}-4x+3=(x-1)(x-3)>0$ sur $D=\intOO{-\infty}{1}\cup\intOO{3}{+\infty}$. $f'(x)=\dfrac{2x-4}{x^{2}-4x+3}$ : $f$ décroît sur
    $\intOO{-\infty}{1}$ et croît sur $\intOO{3}{+\infty}$. Limites : $+\infty$ en $\pm\infty$, $-\infty$ en $1^{-}$ et en $3^{+}$
    (asymptotes $x=1$ et $x=3$). La courbe est symétrique par rapport à la droite $x=2$, car $f(4-x)=f(x)$.
  \item Le système équivaut à $xy=12$ et $\frac xy=3$, d'où $x=3y$, $3y^{2}=12$ et $y=2$ : $\res{(x\pv y)=(6\pv2)}$.
  \item $h(x)=x-\ln(1+x)$ a pour dérivée $\frac{x}{1+x}$, négative sur $\intOO{-1}{0}$ et positive ensuite : $h$ admet un minimum
    $h(0)=0$, donc $\res{\ln(1+x)\leq x}$.
  \item Pour $x>0$, $\ln\left(x^{2}+1\right)=2\ln x+\ln\left(1+\frac{1}{x^{2}}\right)$, donc $\dfrac{\ln\left(x^{2}+1\right)}{x}=\dfrac{2\ln x}{x}+\dfrac{\ln\left(1+\frac{1}{x^{2}}\right)}{x}\to\res{0}$.
  \item Pour $x>0$, $\ln x=ax\iff\varphi(x)=a$ avec $\varphi(x)=\frac{\ln x}{x}$. $\varphi'(x)=\frac{1-\ln x}{x^{2}}$ : $\varphi$ croît de $-\infty$
    à $\frac1e$ sur $\intOF{0}{e}$, puis décroît vers $0$ (en restant positive). Une seule solution si et seulement si
    $\res{a\leq0\text{ ou }a=\frac1e}$ ; deux solutions si $0<a<\frac1e$ ; aucune si $a>\frac1e$.
\end{questions}
\end{corrige}

%% ---------------------------------------------------------------------
\begin{corrige}{10}{Synthèse 2}
\noindent
\begin{minipage}[t]{0.54\linewidth}\vspace{0pt}\raggedright
\begin{questions}
  \item $f'(x)=\frac1x-1=\res{\dfrac{1-x}{x}}$, positive sur $\intOO{0}{1}$, négative sur $\intOO{1}{+\infty}$.
  \item $\lim_{0^{+}}f=-\infty$ ; $f(x)=x\left(\frac{\ln x}{x}-1+\frac2x\right)\to-\infty$ en $+\infty$. Maximum $f(1)=1$.
  \item Sur $\intOF{0}{1}$, $f$ est continue, strictement croissante, de $-\infty$ à $1>0$ : une solution $x_1$. Sur
    $\intFO{1}{+\infty}$, une solution $x_2$ de même. $\res{\text{Deux solutions}}$ : $x_1\approx0{,}159$ et $x_2\approx3{,}146$.
  \item $\ln x\geq x-2\iff f(x)\geq0$ : $\res{S=\intFF{x_1}{x_2}}$.
\end{questions}
\end{minipage}\hfill
\begin{minipage}[t]{0.42\linewidth}\vspace{1mm}
\centering
\begin{tableau}
  \otTab[1.5]{$x$/0.9,$f'(x)$/0.9,$f$/1.9}{$0$,$1$,$+\infty$}
  \tkzTabLine{d,+,z,-,}
  \tkzTabVar{D-/ /$-\infty$,+/$1$,-/$-\infty$}
\end{tableau}
\end{minipage}
\end{corrige}

%% ---------------------------------------------------------------------
\begin{corrige}{11}{Une composée avec le logarithme}
\noindent
\begin{minipage}[t]{0.56\linewidth}\vspace{0pt}\raggedright
\begin{questions}
  \item $x^{2}+1>0$ : le quotient est strictement positif si et seulement si $x^{2}>4$, d'où $D$. $D$ est symétrique par rapport à
    $0$ et $f(-x)=f(x)$ : $f$ est paire.
  \item Sur $D$, $f(x)=\ln\left(x^{2}+1\right)-\ln\left(x^{2}-4\right)$, donc
    $f'(x)=\dfrac{2x}{x^{2}+1}-\dfrac{2x}{x^{2}-4}=\dfrac{2x\left(x^{2}-4-x^{2}-1\right)}{\left(x^{2}+1\right)\left(x^{2}-4\right)}$.
    $f$ est croissante sur $\intOO{-\infty}{-2}$ et décroissante sur $\intOO{2}{+\infty}$.
  \item En $\pm\infty$, le quotient tend vers $1$ : $\res{\lim f=0}$, asymptote $y=0$. En $2^{+}$ et $-2^{-}$, il tend vers $+\infty$ :
    asymptotes $\res{x=2}$ et $\res{x=-2}$.
  \item $x^{2}+1>x^{2}-4>0$ : le quotient est supérieur à $1$, donc $\res{f(x)>0}$.
  \item Avec $X=x^{2}$, $3X^{2}-3X+4$ a un discriminant $9-48<0$ : il est toujours positif. Donc $f''>0$ et
    $\res{f\text{ est convexe}}$ sur chacun des deux intervalles.
\end{questions}
\end{minipage}\hfill
\begin{minipage}[t]{0.4\linewidth}\vspace{0pt}
\centering
\begin{tikzpicture}
\begin{axis}[otSchema, axis lines=middle, width=62mm, height=48mm, xmin=-6.5, xmax=6.5, ymin=-0.3, ymax=3.6,
  xtick={-4,-2,2,4}, ytick={1,2,3}, xlabel={$x$}, xlabel style={at={(1,0.08)},anchor=west},
  grid=major, grid style={otRule!70,line width=0.3pt}]
  \draw[otGreen,dashed,line width=0.6pt] (axis cs:2,-0.3) -- (axis cs:2,3.6);
  \draw[otGreen,dashed,line width=0.6pt] (axis cs:-2,-0.3) -- (axis cs:-2,3.6);
  \addplot[otNavy,line width=1.2pt,domain=2.043:6.5,samples=120] {ln((x^2+1)/(x^2-4))};
  \addplot[otNavy,line width=1.2pt,domain=-6.5:-2.043,samples=120] {ln((x^2+1)/(x^2-4))};
\end{axis}
\end{tikzpicture}\\[0.5mm]
\legende{Courbe paire, convexe, au-dessus de l'axe $y=0$.}
\end{minipage}
\end{corrige}

%% ---------------------------------------------------------------------
\begin{corrige}{12}{Fonction auxiliaire et asymptote oblique}
\partieExo{A}{Une fonction auxiliaire}
\begin{questions}
  \item $g'(x)=2x+\frac1x>0$ : $\res{g\text{ est strictement croissante}}$ sur $\intOO{0}{+\infty}$.
  \item $g(1)=0$ : $\res{g(x)<0\text{ sur }\intOO{0}{1}}$ et $\res{g(x)>0\text{ sur }\intOO{1}{+\infty}}$.
\end{questions}
\partieExo{B}{Étude de la fonction $f$}
\begin{questions}[resume]
  \item En $0^{+}$ : $\frac{\ln x}{x}=\ln x\times\frac1x\to-\infty$, donc $\res{\lim f=+\infty}$. En $+\infty$ : $\frac{\ln x}{x}\to0$, donc $\res{\lim f=+\infty}$.
  \item $f'(x)=1-\dfrac{\frac1x\times x-\ln x}{x^{2}}=1-\dfrac{1-\ln x}{x^{2}}=\dfrac{x^{2}-1+\ln x}{x^{2}}=\res{\dfrac{g(x)}{x^{2}}}$, du signe de $g(x)$.
\end{questions}
\noindent
\begin{minipage}[c]{0.46\linewidth}
\centering
\begin{tableau}
  \otTab[1.6]{$x$/0.9,$f'(x)$/0.9,$f$/1.9}{$0$,$1$,$+\infty$}
  \tkzTabLine{d,-,z,+,}
  \tkzTabVar{D+/ /$+\infty$,-/$0$,+/$+\infty$}
\end{tableau}
\end{minipage}\hfill
\begin{minipage}[c]{0.5\linewidth}
\centering
\begin{tikzpicture}
\begin{axis}[otGraphe, width=64mm, xmin=-0.3, xmax=6.5, ymin=-1.3, ymax=7.5,
  xtick={1,2,3,4,5,6}, ytick={2,4,6}]
  \addplot[otGreen,dashed,line width=0.7pt,domain=0:6.5] {x-1};
  \addplot[otNavy,line width=1.2pt,domain=0.2:6.5,samples=150] {x-1-ln(x)/x};
  \fill[otNavy] (axis cs:1,0) circle (1.4pt);
  \node[anchor=west,font=\scriptsize,text=otGreenDark] at (axis cs:5.1,3.6) {$\Delta$};
  \node[anchor=west,font=\scriptsize,text=otNavy] at (axis cs:0.45,5.6) {$\Cf$};
\end{axis}
\end{tikzpicture}
\end{minipage}
\begin{questions}[start=5]
  \item $f(x)-(x-1)=-\dfrac{\ln x}{x}\to0$ en $+\infty$ : $\res{\Delta\text{ est asymptote à }\Cf}$. Comme $-\frac{\ln x}{x}$ est du signe de
    $-\ln x$, $\Cf$ est \textbf{au-dessus} de $\Delta$ sur $\intOO{0}{1}$ et \textbf{en dessous} sur $\intOO{1}{+\infty}$ ; elles se coupent en $(1\pv0)$.
  \item Voir la figure ci-dessus.
\end{questions}
\end{corrige}

%% ---------------------------------------------------------------------
\begin{corrige}{13}{Injections répétées d'un médicament}
\noindent
\begin{minipage}[t]{0.54\linewidth}\vspace{0pt}\raggedright
\partieExo{A}{Une seule injection}
\begin{questions}
  \item Il reste 70\,\% chaque heure : $\res{u_n=20\times0{,}7^{n}}$.
  \item $0{,}7^{n}<0{,}05\iff n\ln0{,}7<\ln0{,}05\iff n>\dfrac{\ln0{,}05}{\ln0{,}7}\approx8{,}4$, car $\ln0{,}7<0$ :
    au bout de $\res{9\text{ heures}}$.
\end{questions}
\partieExo{B}{Injections toutes les heures}
\begin{questions}[resume]
  \item Il reste $0{,}7\,v_n$, auxquels on ajoute $5$~mg. Puis $w_{n+1}=0{,}7\,v_n+5-\frac{50}{3}=0{,}7\left(v_n-\frac{50}{3}\right)=0{,}7\,w_n$ :
    $(w_n)$ est géométrique de raison $0{,}7$ et $w_0=-\frac{35}{3}$.
  \item $v_n=\frac{50}{3}+w_n=\res{\frac{50}{3}-\frac{35}{3}\times0{,}7^{n}}$ et $\res{\lim v_n=\frac{50}{3}\approx16{,}7}$ : la quantité se stabilise.
  \item $v_n\geq15\iff\frac{35}{3}\times0{,}7^{n}\leq\frac53\iff0{,}7^{n}\leq\frac17\iff n\geq\frac{\ln(1/7)}{\ln0{,}7}\approx5{,}46$ :
    $\res{6\text{ réinjections}}$ ($v_5\approx14{,}71$ et $v_6\approx15{,}29$).
\end{questions}
\end{minipage}\hfill
\begin{minipage}[t]{0.42\linewidth}\vspace{0pt}
\centering
\begin{tikzpicture}
\begin{axis}[otSchema, width=66mm, height=54mm, xmin=-0.8, xmax=14.8, ymin=0, ymax=19,
  xtick={0,2,...,14}, ytick={5,10,15}, xlabel={$n$}, xlabel style={at={(1,0)},anchor=west},
  grid=major, grid style={otRule!70,line width=0.3pt}]
  \addplot[ybar,bar width=2.2mm,fill=otNavySoft,draw=otNavy,line width=0.4pt,samples at={0,...,14}]
    {50/3-35/3*0.7^x};
  \draw[otNavy,dashed,line width=0.6pt] (axis cs:-0.8,16.667) -- (axis cs:14.8,16.667);
  \draw[otGreen,line width=0.7pt] (axis cs:-0.8,15) -- (axis cs:14.8,15);
  \node[anchor=south east,font=\scriptsize,text=otGreenDark] at (axis cs:14.8,12.6) {seuil $15$};
\end{axis}
\end{tikzpicture}\\[0.5mm]
\legende{Quantité après chaque réinjection ; en pointillés, la limite $\frac{50}{3}$.}
\end{minipage}
\end{corrige}

%% ---------------------------------------------------------------------
\begin{corrige}{14}{La série harmonique diverge}
\noindent
\begin{minipage}[t]{0.56\linewidth}\vspace{0pt}\raggedright
\partieExo{A}{Une inégalité}
\begin{questions}
  \item $\varphi(x)=x-1-\ln x$ : $\varphi'(x)=1-\frac1x=\frac{x-1}{x}$, d'où un minimum $\varphi(1)=0$. Donc $\varphi\geq0$ : $\res{\ln x\leq x-1}$.
  \item Avec $x=1+\frac1k>0$ : $\res{\ln\left(1+\frac1k\right)\leq\frac1k}$.
\end{questions}
\partieExo{B}{Divergence de $(H_n)$}
\begin{questions}[resume]
  \item $\ln\left(\frac{k+1}{k}\right)=\ln(k+1)-\ln k$. La somme est télescopique :
    $\big(\ln2-\ln1\big)+\big(\ln3-\ln2\big)+\dots+\big(\ln(n+1)-\ln n\big)=\res{\ln(n+1)}$.
  \item En sommant les inégalités de la question 2 : $\res{\ln(n+1)\leq H_n}$. Comme $\lim\ln(n+1)=+\infty$, par comparaison
    $\res{\lim H_n=+\infty}$.
  \item Il suffit que $\ln(n+1)>10$, soit $n>e^{10}-1\approx22\,025{,}5$ : pour $\res{n\geq22\,026}$, $H_n>10$.
  \item La fonction renvoie $\res{83}$ : il faut $83$ termes pour dépasser $5$. La suite tend vers $+\infty$ très lentement,
    à peu près comme $\ln n$ : on montre que $H_n-\ln n$ converge vers $\gamma\approx0{,}577$, la constante d'Euler.
\end{questions}
\end{minipage}\hfill
\begin{minipage}[t]{0.4\linewidth}\vspace{0pt}
\centering
\begin{tikzpicture}
\begin{axis}[otSchema, width=62mm, height=52mm, xmin=0, xmax=21, ymin=0, ymax=4,
  xtick={5,10,15,20}, ytick={1,2,3}, xlabel={$n$}, xlabel style={at={(1,0)},anchor=west},
  grid=major, grid style={otRule!70,line width=0.3pt}]
  \addplot[otGreen,line width=1pt,domain=1:20.5,samples=80] {ln(x+1)};
  \addplot[only marks,mark=*,mark size=1.3pt,otNavy] coordinates {
    (1,1) (2,1.5) (3,1.8333) (4,2.0833) (5,2.2833) (6,2.45) (7,2.5929) (8,2.7179) (9,2.829) (10,2.929)
    (11,3.0199) (12,3.1032) (13,3.1801) (14,3.2516) (15,3.3182) (16,3.3807) (17,3.4396) (18,3.4951) (19,3.5477) (20,3.5977)};
  \node[anchor=north west,font=\scriptsize,text=otGreenDark] at (axis cs:12,2.5) {$\ln(n+1)$};
  \node[anchor=south east,font=\scriptsize,text=otNavy] at (axis cs:9,3.1) {$H_n$};
\end{axis}
\end{tikzpicture}\\[0.5mm]
\legende{Les points $H_n$ restent au-dessus de la courbe de $\ln(n+1)$.}
\end{minipage}
\end{corrige}
